\section{Från buffer till ström}
\label{sec:frombuffer}

I \kapref{ch:frames} antog vi att en komplett frame låg i en buffer. Det gör den aldrig.

En seriell länk levererar en byte i taget, när den kommer, och en nod som väntar på nästa byte får
inte stanna och vänta --- den har annat att göra. Mottagningen måste alltså vara
\textbf{inkrementell}: ta emot en byte, gör vad som går att göra med den, och återvänd.

Det är precis vad en tillståndsmaskin är till för.

\subsection{Tillstånden}

Ett tillstånd per fält, plus ett som betyder klar:

\begin{console}
WaitForSof1 -> WaitForSof2 -> WaitForLen  -> WaitForType -> WaitForDst
            -> WaitForSrc  -> WaitForSeq1 -> WaitForSeq2 -> WaitForPayload
            -> WaitForChk1 -> WaitForChk2 -> Ready
\end{console}

Parserns publika gränssnitt är litet:

\begin{cppcode}
class Parser final
{
public:
    /** Feed the parser one byte. Returns false on a parse error. */
    bool processByte(std::uint8_t byte) noexcept;

    /** True when a complete frame has been framed. */
    bool isFrameReady() const noexcept;

    /** Validate and copy out the framed data. */
    bool extractFrame(Frame& frame) const noexcept;

    /** Discard any partial frame and start over. */
    void reset() noexcept;
};
\end{cppcode}

\subsection{Två ansvar, medvetet åtskilda}

Lägg märke till att \code{isFrameReady()} och \code{extractFrame()} är två olika saker, och att
det är avsiktligt:

\begin{itemize}
  \item \code{isFrameReady()} säger att \textbf{framningen} är klar --- rätt antal bytes har
    kommit in, på rätt platser.
  \item \code{extractFrame()} säger att framen är \textbf{giltig} --- SOF stämmer, längden är
    rimlig, typen finns, och checksumman går ihop.
\end{itemize}

En frame kan vara färdigframad och ändå trasig. Det är själva normalfallet när något gått fel på
ledningen, och därför måste de två frågorna ställas var för sig.

\begin{cppcode}
bool Parser::extractFrame(Frame& frame) const noexcept
{
    if (!isFrameReady()) { return false; }

    // Deserialize validates SOF, payload length, frame type and checksum, so a frame that was
    // framed correctly but arrived corrupted is rejected here.
    return frame.deserialize(myBuf, static_cast<std::size_t>(myParsedBytes));
}
\end{cppcode}

Valideringen ligger alltså i \code{Frame::deserialize()} och inte i tillståndsmaskinen. Det är
inte en slump: tillståndsmaskinen vet var fälten \emph{ligger}, framen vet vad de \emph{betyder},
och att blanda ihop de två ansvaren ger en parser som måste ändras varje gång protokollet gör det.

\begin{aside}
\note \code{extractFrame()} är \code{const} och kan därför inte nollställa parsern. Det är ett
medvetet val: den som anropar den vet om den vill försöka igen eller kasta bort framen, och gör
\code{reset()} själv. Samma mönster återkommer i drivrutinen: en metod rapporterar, anroparen
bestämmer.
\end{aside}
